\section{Reduced Artin--Schreier parameters and a common pole bound}

Let $p$ be prime, put $k=\mathbf F_p$, and write
$\wp(x)=x^p-x$. In any field containing $k$, this is a
$k$-linear map: Frobenius is additive in characteristic $p$,
and $c^p=c$ for $c\in k$. Thus its image is a $k$-subspace
of the additive group, and the quotient by its image is a
$k$-vector space.

We give proofs of the reduced representatives in
\cite{Camina1997}, Lemmas~2--3, and of the finite-kernel
argument in the proof of her Theorem~4. The latter is expressed
directly through additive characters instead of cohomology.
The final parameter choice uses a single bound for all the
finitely many desired pole inequalities. Over Laurent fields this
choice also supplies irreducibility without a Frattini hypothesis:
a sufficiently deep reduced pole cannot be an Artin--Schreier image.

For the first three lemmas, let $K=k((T))$ and let $v_T$ be
its additive order valuation, with $v_T(T)=1$ and
$v_T(0)=+\infty$.

\begin{lemma}[Positive-order series are Artin--Schreier images]
\label{lem:artin-schreier-positive-tail}
\lean{LaurentSeries.summable_pow_prime_pow, LaurentSeries.pos_orderTop_tsum_prime_pow, LaurentSeries.artinSchreier_neg_tsum_prime_pow, LaurentSeries.exists_pos_orderTop_pow_sub_eq}
\leanok
For every $f\in T k[[T]]$, the series
\[
 h=-\sum_{j=0}^{\infty} f^{p^j}
\]
converges in $T k[[T]]$ and satisfies $\wp(h)=f$.
\end{lemma}

\begin{proof}
If $f=0$, take $h=0$. Otherwise $v_T(f)\ge1$ and
$v_T(f^{p^j})=p^jv_T(f)$ tends to infinity.
The partial sums are Cauchy and converge by completeness of
$k((T))$. Each partial sum has order at least one, as does
the limit: for any convergent sequence of elements of order
at least one, a sufficiently small difference between a term
and its limit gives the same bound on the limit.
For $h_N=-\sum_{j=0}^{N-1}f^{p^j}$, additivity of Frobenius
gives the telescoping identity
\[
 \wp(h_N)
   =-\sum_{j=1}^{N}f^{p^j}+\sum_{j=0}^{N-1}f^{p^j}
   =f-f^{p^N}.
\]
The last term tends to $f$. Addition and multiplication are
continuous for the valuation topology, so $\wp(h_N)$ also
tends to $\wp(h)$. Uniqueness of the limit gives
$\wp(h)=f$.
\end{proof}

\begin{lemma}[Unique reduced representative]
\label{lem:artin-schreier-reduced-representative}
\uses{lem:artin-schreier-positive-tail,lem:artin-schreier-no-root-from-pole}
\source{camina1997:page-0005}
\lean{Camina1997.ArtinSchreier.existsUnique_reduced, Camina1997.ArtinSchreier.IsReduced.finite_support, Camina1997.ArtinSchreier.IsReduced.eq_single_or_exists_pole, Camina1997.ArtinSchreier.exists_reduced_of_not_image}
\leanok
Every class in $K/\wp(K)$ has a unique representative of the
form
\[
 r=c+\sum_{\substack{m>0\\p\nmid m}}a_m T^{-m},
 \qquad c,a_m\in k,
\]
where only finitely many $a_m$ are nonzero.
In particular every $\gamma\in K$ can be written
$\gamma=r+\wp(\mu)$ with $\mu\in K$ and $r$ reduced in
this sense. If $\gamma\notin\wp(K)$, then $r\ne0$,
$v_T(r)\le0$, and a negative value of $v_T(r)$ is not
divisible by $p$.
\end{lemma}

\begin{proof}
Split a Laurent series into its constant coefficient, its
finite negative-degree part, and its positive-order part.
The positive-order part belongs to $\wp(K)$ by the preceding
lemma. For a negative monomial $aT^{-n}$ with $a\in k$ and
$n>0$, write $n=p^e m$ with $p\nmid m$. Repeated division
of $n$ by $p$ produces such $e,m$, since the positive integer
strictly decreases at each division. We have
\[
 \wp\left(\sum_{j=0}^{e-1}aT^{-p^j m}\right)
   =aT^{-p^e m}-aT^{-m},
\]
where the sum is zero if $e=0$.
Indeed the two sums in its Artin--Schreier image telescope,
using $a^p=a$. Thus $aT^{-n}$ is congruent modulo
$\wp(K)$ to $aT^{-m}$. Apply this to each of the finitely
many negative monomials and collect the coefficients with
equal $m$. Together with the positive-tail solution, this
gives the stated expression $\gamma=r+\wp(\mu)$.

For uniqueness, suppose the difference of two reduced
expressions is $\wp(x)$. Their difference is reduced as
well. If it has a nonzero negative term, choose its largest
pole exponent $m$. Its valuation is $-m$, with $p\nmid m$,
contradicting the pole obstruction lemma.
The difference is therefore a constant $c\in k$.
If $v_T(x)<0$, then $p v_T(x)<v_T(x)$, so
$v_T(x^p-x)=p v_T(x)<0$, which cannot be the valuation of
the constant $c$. Hence $x\in k[[T]]$.
Taking constant coefficients in $x^p-x=c$ gives
$c=x_0^p-x_0=0$, since $x_0\in k$.
The reduced expressions are equal, and uniqueness of Laurent
coefficients gives equality of all their displayed
coefficients. Finally a reduced nonzero expression is either
a nonzero constant or has lowest exponent $-m$ with
$m>0$ and $p\nmid m$, proving the additional assertions.
\end{proof}

\begin{lemma}[Basis of the Artin--Schreier quotient]
\label{lem:artin-schreier-quotient-basis}
\uses{lem:artin-schreier-reduced-representative}
\source{camina1997:page-0005}
\lean{Camina1997.ArtinSchreier.reducedQuotientEquiv, Camina1997.ArtinSchreier.quotientBasis, Camina1997.ArtinSchreier.quotientBasis_apply, Camina1997.ArtinSchreier.not_moduleFinite_quotient}
\leanok
The classes of
\[
 \{1\}\ \cup\ \{T^{-m}:m>0,\ p\nmid m\}
\]
form a $k$-basis of $K/\wp(K)$. In particular this quotient
has infinite dimension over $k$.
\end{lemma}

\begin{proof}
Existence of a reduced representative says exactly that these
classes span the quotient. A finite linear combination which
represents zero is a reduced representative of the zero
class. Uniqueness forces all of its coefficients to be zero,
so the displayed classes are linearly independent.
The exponents $1+np$, for $n\ge0$, are distinct positive
integers not divisible by $p$. They give an infinite subset
of the basis, proving infinite dimension.
\end{proof}

For a field inclusion $F\subseteq K$ of characteristic $p$,
write
\[
 \operatorname{res}:F/\wp(F)\longrightarrow K/\wp(K)
\]
for the $k$-linear map induced by the inclusion. It is
well-defined because the inclusion commutes with $\wp$.

\begin{lemma}[The restriction kernel is finite]
\label{lem:artin-schreier-restriction-finite-kernel}
\source{camina1997:page-0009}
\lean{Camina1997.ArtinSchreier.restriction, Camina1997.ArtinSchreier.exists_root_of_mem_ker_restriction, Camina1997.ArtinSchreier.finite_wpKernel, Camina1997.ArtinSchreier.finite_ker_restriction, Camina1997.ArtinSchreier.moduleFinite_ker_restriction}
\leanok
For a finite Galois extension $K/F$, the kernel of the map
\[
 F/\wp(F)\longrightarrow K/\wp(K)
\]
is a finite set, hence a finite-dimensional $\mathbf F_p$-vector space.
\end{lemma}

\begin{proof}
For every kernel class $q$, choose a representative $a_q\in F$
and a root $\theta_q\in K$ with $\wp(\theta_q)=a_q$.
For $g\in G=\operatorname{Gal}(K/F)$, additivity of $\wp$
and the fact that $g$ fixes $a_q$ imply
$\wp(g(\theta_q)-\theta_q)=0$.
The root set $W=\{x\in K:x^p-x=0\}$ is finite because the
nonzero polynomial $X^p-X$ has at most $p$ roots.
Thus $q\mapsto(g\mapsto g(\theta_q)-\theta_q)$ takes values
in the finite function set $W^G$.
If two classes give the same function, $\theta_q-\theta_r$
is fixed by every element of $G$, and hence equals some $b\in F$
by the Galois fixed-field theorem. Then
$a_q-a_r=\wp(\theta_q-\theta_r)=\wp(b)$, so $q=r$.
The displayed map is injective. Consequently its domain is finite;
its finitely many elements span it as a vector space.
\end{proof}

\begin{lemma}[Restriction from a Laurent field has infinite-dimensional image]
\label{lem:artin-schreier-restriction-infinite-image}
\uses{lem:artin-schreier-quotient-basis,lem:artin-schreier-restriction-finite-kernel}
\source{camina1997:page-0009}
\lean{Camina1997.ArtinSchreier.not_moduleFinite_range_restriction}
\leanok
Let $F=k((t))$ and let $K/F$ be any finite Galois extension.
The image of
$\operatorname{res}:F/\wp(F)\to K/\wp(K)$ has infinite
dimension over $k$.
\end{lemma}

\begin{proof}
The domain is infinite-dimensional by the Laurent quotient
basis, and the kernel is finite-dimensional by the preceding
lemma. If the image had a finite basis, choose preimages of
its finitely many basis vectors and a finite basis of the
kernel. These two lists would span the domain: subtract from
any vector the linear combination of the chosen preimages
having the same image, and the difference lies in the kernel.
That would make the domain finite-dimensional, a contradiction.
\end{proof}

\begin{lemma}[Infinite-dimensional subspaces have unbounded reduced poles]
\label{lem:artin-schreier-unbounded-reduced-poles}
\uses{lem:artin-schreier-reduced-representative,lem:artin-schreier-quotient-basis}
\source{camina1997:page-0009}
\lean{Camina1997.ArtinSchreier.finite_set_reduced_orderTop_ge, Camina1997.ArtinSchreier.exists_reduced_pole_lt_of_not_moduleFinite, Camina1997.ArtinSchreier.exists_reduced_pole_lt_in_coset}
\leanok
Let $K=k((T))$ and let $U$ be an infinite-dimensional
$k$-subspace of $K/\wp(K)$.
For every integer $B\ge0$ and every class $q$, there is a class
in $q+U$ whose reduced representative $r$ satisfies $v_T(r)<-B$.
\end{lemma}

\begin{proof}
A reduced representative of valuation at least $-B$ has no
nonzero coefficients outside degrees $-B,\ldots,0$.
Its coefficient tuple in these degrees determines it, and
there are finitely many such tuples because $k=\mathbf F_p$
is finite. Thus bounded reduced representatives form a finite set.
If every class of $q+U$ had a bounded reduced representative,
the unique reduced representative of $q+u$ would give an injection
from $U$ into this finite set. Then $U$ would be finite and its
finitely many elements would span it, contradicting infinite dimension.
\end{proof}

\begin{lemma}[A reduced parameter below any pole bound]
\label{lem:artin-schreier-deep-parameter-adjustment}
\uses{lem:artin-schreier-restriction-infinite-image,lem:artin-schreier-unbounded-reduced-poles,lem:artin-schreier-reduced-representative}
\source{camina1997:page-0009}
\lean{Camina1997.ArtinSchreier.exists_reduced_pole_adjustment, Camina1997.ArtinSchreier.exists_reduced_pole_adjustment_order}
\leanok
Suppose $F\subseteq K=k((T))$ is finite Galois and $F/\wp(F)$
is infinite-dimensional over $k$; in particular this holds when
$F$ is a Laurent-series field over $k$. For every $\gamma\in K$
and integer $B\ge0$, there exist $b\in F$, $\mu\in K$ and a reduced
$r\in K$ with
\[
 r=\gamma+b-\wp(\mu),\qquad v_T(r)=-m,
 \qquad m>B,\quad p\nmid m.
\]
\end{lemma}

\begin{proof}
The restriction image $U$ is infinite-dimensional. Apply the
preceding lemma to the affine coset $[\gamma]+U$ to obtain a
reduced $r$ with $v_T(r)<-B$. Choose $b\in F$ whose restricted
class is $[r]-[\gamma]$. Then $[\gamma+b-r]=0$ in $K/\wp(K)$,
so $\gamma+b-r=\wp(\mu)$ for some $\mu\in K$.
The negative leading order of a reduced representative is $-m$
with $p\nmid m$, and the bound on its order gives $m>B$.
\end{proof}

\begin{theorem}[One reduced parameter below a prescribed pole bound]
\label{thm:witt-simultaneous-pole-adjustment}
\uses{lem:artin-schreier-deep-parameter-adjustment}
\source{camina1997:page-0005,camina1997:page-0009,witt1936:page-0004}
\lean{Camina1997.Witt.adjustedDelta, Camina1997.Witt.adjustedDelta_one, Camina1997.Witt.adjustedDelta_twoCocycle, Camina1997.Witt.adjustedDelta_parameter, Camina1997.Witt.exists_adjusted_witt_equations}
\leanok
Suppose $F\subseteq K=k((T))$ is finite Galois,
with group $G$, and $F/\wp(F)$ has infinite dimension over $k$.
This hypothesis holds in particular if $F\simeq k((t))$. Let $c:G\times G\to k$ and let
$\delta_s,\gamma\in K$ satisfy
\[
 \delta_1=0,\qquad
 \delta_s+s(\delta_t)-\delta_{st}=c(s,t),\qquad
 s(\gamma)-\gamma=\wp(\delta_s).
\]
For every integer $B\ge0$, there exist $b\in F$ and
$\mu\in K$ such that
\[
 \widehat\gamma=\gamma+b-\wp(\mu),\qquad
 \widehat\delta_s=\delta_s-(s(\mu)-\mu)
\]
satisfy the same three equations with $\gamma,\delta$
replaced by $\widehat\gamma,\widehat\delta$ and with the
same $c$. Moreover $\widehat\gamma$ is reduced and
\[
 v_T(\widehat\gamma)=-m<-B
 \quad\text{for some }m>0\text{ with }p\nmid m.
\]
\end{theorem}

\begin{proof}
Apply Lemma~\ref{lem:artin-schreier-deep-parameter-adjustment}
to $\gamma$ and $B$. It supplies $b\in F$, $\mu\in K$, and a
reduced $r=\gamma+b-\wp(\mu)$ of order $-m$, where $m>B$
and $p\nmid m$. Put $\widehat\gamma=r$ and define
$\widehat\delta_s=\delta_s-(s(\mu)-\mu)$.

Since $b\in F$, every $s\in G$ fixes $b$.
Additivity of $\wp$ and its commutation with automorphisms give
\[
\begin{aligned}
 s(\widehat\gamma)-\widehat\gamma
 &=s(\gamma)-\gamma-\wp(s(\mu)-\mu)\\
 &=\wp(\delta_s-(s(\mu)-\mu))
  =\wp(\widehat\delta_s).
\end{aligned}
\]
Also $\widehat\delta_1=\delta_1=0$.
For the other equation, the correction terms cancel because
\[
 (s(\mu)-\mu)+s(t(\mu)-\mu)-(st(\mu)-\mu)=0.
\]
Thus
$\widehat\delta_s+s(\widehat\delta_t)-\widehat\delta_{st}
 =\delta_s+s(\delta_t)-\delta_{st}=c(s,t)$,
as required.
\end{proof}

Let $d_1,\ldots,d_r$ be the depths of finitely many nonidentity
elements of the quotient. Choose
\[
 B=\max\{0,d_1,\ldots,d_r\}.
\]
The same parameter then satisfies
$v_T(\widehat\gamma)<-d_i$ for every $i$.
This establishes the simultaneous pole condition used in
Camina's Theorem~4; the comparison of the depths upstairs and
downstairs requires the subsequent ramification arguments.


\begin{theorem}[A central lift with an arbitrarily deep reduced parameter]
\label{thm:witt-reduced-central-lift}
\uses{lem:central-extension-factor-set,thm:witt-additive-equations,thm:witt-simultaneous-pole-adjustment,lem:artin-schreier-no-root-from-pole,lem:artin-schreier-irreducible,lem:witt-semilinear-quotient-lifts,lem:full-galois-action-from-degree}
\source{camina1997:page-0009,witt1936:page-0004}
\lean{Camina1997.Witt.exists_reduced_galois_centralLift, Camina1997.Witt.not_moduleFinite_quotient_of_laurentEquiv, Camina1997.Witt.exists_reduced_galois_centralLift_laurentBase}
\leanok
Let $F\subseteq K=k((T))$ be finite Galois, let
$G=\operatorname{Gal}(K/F)$, and suppose $F/\wp(F)$ has
infinite dimension over $k$, as it does when $F\simeq k((t))$.
Let
\[
 1\longrightarrow(k,+)\xrightarrow{\iota}E
 \xrightarrow{q}G\longrightarrow1
\]
be a central extension with $E$ finite. For every integer $B\ge0$
there are a reduced $r\in K$ and an integer $m>B$ with $p\nmid m$
and $v_T(r)=-m$ such that
\[
 L=K[Y]/(Y^p-Y-r)
\]
is a field of degree $p$ over $K$, is Galois over $F$, and admits
an isomorphism $\rho:E\simeq\operatorname{Gal}(L/F)$ whose
restriction to $K$ is $q$.
No inclusion of the kernel in $\Phi(E)$ is required.
\end{theorem}

\begin{proof}
Choose a normalized section of $q$ and its central factor set $c$.
Witt's additive equations give $\delta_s,\gamma\in K$ with
\[
 \delta_1=0,\qquad
 \delta_s+s(\delta_t)-\delta_{st}=c(s,t),\qquad
 s(\gamma)-\gamma=\wp(\delta_s).
\]
The preceding simultaneous adjustment gives $r$ of order $-m$,
with $m>B$ and $p\nmid m$, and corrected increments satisfying
the same equations for $c$. The pole obstruction shows that
$r\notin\wp(K)$, so $Y^p-Y-r$ is irreducible. Thus its quotient
$L$ is a field and has the basis $1,y,\ldots,y^{p-1}$ over $K$.
The semilinear quotient automorphisms, with these corrected
increments, give a faithful action of $E$ on $L$ fixing $F$;
their restrictions to $K$ are the prescribed elements $q(e)$.
As a set $E$ is the product of $G$ and its kernel, hence
\[
 |E|=p|G|=p[K:F]=[L:F].
\]
Lemma~\ref{lem:full-galois-action-from-degree} proves that the
extension is Galois and that the action is its full Galois group.
If $F$ is identified with $k((t))$, its field isomorphism induces
an isomorphism of Artin--Schreier quotients, because it commutes
with $\wp$ and fixes the prime field. The quotient basis then
proves the infinite-dimensional hypothesis.

This argument combines Witt's explicit lifted action with the
parameter adjustment in Camina's proof of Theorem~4. Here the
negative reduced pole proves properness directly, in place of
the Frattini argument used for an arbitrary characteristic-$p$
base field.
\end{proof}

\begin{lemma}[Semilinear lifts preserve the weighted valuation]
\label{lem:artin-schreier-semilinear-valuation}
\uses{lem:artin-schreier-weighted-valuation}
\lean{Camina1997.ArtinSchreier.weightedOrder_semilinear, Camina1997.ArtinSchreier.uniformContinuous_semilinear}
\leanok
Let $K=k((T))$, let $r\in K$ satisfy $v_T(r)=-m$ with
$m>0$ and $p\nmid m$, and put $L=K[Y]/(Y^p-Y-r)$.
Write $y$ for the class of $Y$ and $w$ for the weighted valuation
\[
 w\left(\sum_{i=0}^{p-1}a_i y^i\right)
   =\min_{0\le i<p}\bigl(pv_T(a_i)-mi\bigr).
\]
Suppose $s$ is a ring automorphism of $K$ preserving $v_T$,
and $e$ is a ring automorphism of $L$ whose restriction to $K$
is $s$. Then $w(e(x))=w(x)$ for every $x\in L$, and $e$
is uniformly continuous for the topology defined by $w$.
\end{lemma}

\begin{proof}
Put $z=e(y)$. Applying $e$ to the defining equation of $y$
gives $z^p-z=s(r)$. Since $w|_K=pv_T$, its right-hand side
has value $-pm<0$. Thus $z\ne0$ and $w(z)$ is a finite integer.
If $w(z)\ge0$, the valuation inequality would give
$w(z^p-z)\ge\min\{pw(z),w(z)\}\ge0$, a contradiction.
Consequently $w(z)<0$. Since $p>1$, the two summands
$z^p$ and $-z$ have distinct values $pw(z)<w(z)$;
therefore $w(z^p-z)=pw(z)$. It follows that $w(z)=-m$.

Write $x=\sum_{i=0}^{p-1}a_i y^i$. Then
$e(x)=\sum_{i=0}^{p-1}s(a_i)z^i$, and each term has value
\[
 w(s(a_i)z^i)=pv_T(s(a_i))+i w(z)=pv_T(a_i)-mi.
\]
Each displayed value is at least $w(x)$, so the valuation
inequality for a finite sum gives $w(e(x))\ge w(x)$.
The inverse maps satisfy the same hypotheses. Indeed,
$e^{-1}(a)=s^{-1}(a)$ on $K$, and applying valuation preservation
by $s$ to $s^{-1}(a)$ gives $v_T(s^{-1}(a))=v_T(a)$.
Apply the inequality just proved to $e^{-1}$ and $e(x)$:
\[
 w(x)=w(e^{-1}(e(x)))\ge w(e(x)).
\]
The two inequalities give the required equality, also when
$x=0$ and the values are $+\infty$.
Finally $e(x)-e(x')=e(x-x')$, so the equality just proved
preserves each condition $w(x-x')>n$. These conditions are
a basis of entourages for the valuation uniformity, which
proves uniform continuity.
\end{proof}


\begin{theorem}[Central lifting in Laurent coordinates]
\label{thm:witt-laurent-central-lift}
\uses{thm:witt-reduced-central-lift,lem:artin-schreier-semilinear-valuation,lem:artin-schreier-weighted-valuation,lem:artin-schreier-laurent-parameter}
\source{camina1997:page-0005,camina1997:page-0009}
\lean{Camina1997.Witt.exists_laurent_centralLift}
\leanok
Put $K=k((T))$ and write $v$ for its order valuation.
Let $F$ be a field with an isomorphism $e_F:F\simeq K$, and put
$v_F=v\circ e_F$. Let $\alpha:F\hookrightarrow K$ be a field
embedding making $K/F$ finite Galois, and write
$G=\operatorname{Gal}_{\alpha}(K/F)$ for its Galois group.
Suppose
\[
 v(g(x))=v(x)\quad(g\in G,\ x\in K),\qquad
 v(\alpha(a))=|G|v_F(a)\quad(a\in F).
\]
Let $E$ be finite and let $q:E\twoheadrightarrow G$ have a
central kernel of order $p$. There is a field embedding
$\beta:F\hookrightarrow K$ making $K/F$ finite Galois and an
isomorphism $\rho:E\simeq\operatorname{Gal}_{\beta}(K/F)$ such that
\[
 v(\rho(e)(x))=v(x)\quad(e\in E,\ x\in K),\qquad
 v(\beta(a))=|E|v_F(a)\quad(a\in F).
\]
Here $\alpha$ and $\beta$ are possibly different embeddings into
the same Laurent-series field.
\end{theorem}

\begin{proof}
A group of order $p$ is cyclic: the order of any nonidentity
element divides $p$ and is greater than one, hence equals $p$.
An isomorphism from $(k,+)$ onto $\ker q$ therefore presents
$q$ as a central extension of the form used in
Theorem~\ref{thm:witt-reduced-central-lift}.
Apply that theorem with $B=0$. It gives a reduced parameter
$r\in K$ with $v(r)=-m$, $m>0$, and $p\nmid m$, a degree-$p$
field
\[
 A=K[Y]/(Y^p-Y-r),
\]
and an isomorphism $\tau:E\simeq\operatorname{Gal}(A/F)$ whose
restriction to $K$ is $q$. Write $j:K\hookrightarrow A$ for
the coefficient inclusion and $w$ for the weighted valuation.
Every $q(e)$ preserves $v$ by hypothesis, so
Lemma~\ref{lem:artin-schreier-semilinear-valuation} gives
$w(\tau(e)(z))=w(z)$ for every $z\in A$.

The weighted valuation is surjective, so choose $\pi\in A$
with $w(\pi)=1$. The Laurent parameter theorem constructs a
$k$-algebra isomorphism $u:K\simeq A$ with
$w(u(x))=v(x)$. Define
\[
 \beta=u^{-1}\circ j\circ\alpha,
 \qquad
 \rho(e)=u^{-1}\circ\tau(e)\circ u.
\]
The first map is a field embedding, and $u$ identifies the
extension of $F$ through $\beta$ with $A/F$. It is therefore
finite Galois; conjugation by $u$ identifies its full Galois
group with that of $A/F$. This proves the assertion about
$\rho$.
For $e\in E$ and $x\in K$,
\[
 v(\rho(e)(x))=w(\tau(e)(u(x)))=w(u(x))=v(x).
\]
Also $w(j(b))=pv(b)$ for $b\in K$, so
\[
 v(\beta(a))=w(j(\alpha(a)))=p v(\alpha(a))
   =p|G|v_F(a)=|E|v_F(a),
\]
where the last equality follows from $|\ker q|=p$ and
surjectivity of $q$.

The parameter adjustment used here is the one from Camina's
proof of Theorem~4. It supplies a negative reduced pole at every
central step, without requiring a Frattini kernel; the Laurent
coordinate construction gives explicitly the field identification
used in the proof of her Theorem~3.
\end{proof}


\begin{theorem}[Finite $p$-groups as Galois groups in Laurent coordinates]
\label{thm:witt-finite-laurent-realization}
\uses{lem:finite-p-group-central-prime-step,thm:witt-laurent-central-lift}
\source{camina1997:page-0005}
\lean{Camina1997.Witt.exists_laurent_galois_realization}
\leanok
Put $K=k((T))$ with its order valuation $v$.
Let $F$ be a field equipped with an isomorphism $e_F:F\simeq K$,
and set $v_F=v\circ e_F$. For every finite $p$-group $H$ there
are a field embedding $\beta:F\hookrightarrow K$ making $K/F$
finite Galois and an isomorphism
$\rho:H\simeq\operatorname{Gal}_{\beta}(K/F)$ such that
\[
 v(\rho(h)(x))=v(x)\quad(h\in H,\ x\in K),\qquad
 v(\beta(a))=|H|v_F(a)\quad(a\in F).
\]
\end{theorem}

\begin{proof}
We use strong induction on $|H|$. If $H$ is trivial, take
$\beta=e_F$. Then $\beta$ is onto, so $K/F$ has degree one,
its Galois group is trivial, and its unique action preserves $v$.
The second identity reduces to the definition of $v_F$.

Suppose $H$ is nontrivial. Choose a central subgroup $N$ of
order $p$. It is normal, and $Q=H/N$ is a $p$-group of strictly
smaller order. The induction hypothesis gives an embedding
$\alpha:F\hookrightarrow K$ and an isomorphism
$\sigma:Q\simeq\operatorname{Gal}_{\alpha}(K/F)$ with the two
valuation properties. Compose the quotient map $H\to Q$ with
$\sigma$ to obtain
\[
 q:H\twoheadrightarrow\operatorname{Gal}_{\alpha}(K/F).
\]
Because $\sigma$ is injective, $\ker q=N$, so this kernel is
central of order $p$. Every automorphism in the target of $q$
is $\sigma(s)$ for some $s\in Q$ and hence preserves $v$.
Moreover $|Q|=|\operatorname{Gal}_{\alpha}(K/F)|$, so the
inductive base-valuation identity has exactly the hypothesis
required in Theorem~\ref{thm:witt-laurent-central-lift}.
That theorem produces the desired embedding $\beta$ and full
Galois-group isomorphism for $H$, with both valuation identities.
This completes the induction.

This construction implements the realization used in Camina's
Theorem~3 directly in Laurent coordinates. Choosing a deep reduced
parameter at each central step allows all finite $p$-groups to be
handled by the same induction, including elementary abelian groups.
\end{proof}

\begin{lemma}[Valuation-preserving Laurent actions restrict faithfully]
\label{lem:laurent-action-nottingham-restriction}
\uses{lem:nottingham-laurent-extension,lem:pgroup-power-series-normalization}
\source{camina1997:page-0005}
\lean{Camina1997.Nottingham.restrictLaurentAction, Camina1997.Nottingham.coe_restrictLaurentAction, Camina1997.Nottingham.restrictLaurentAction_injective_iff, Camina1997.Nottingham.ofLaurentAction, Camina1997.Nottingham.toLaurent_ofLaurentAction, Camina1997.Nottingham.ofLaurentAction_injective_iff, Camina1997.Nottingham.isClosedEmbedding_ofLaurentAction}
\leanok
Let $G$ be a group acting on $k((T))$ by $k$-algebra automorphisms,
and suppose each automorphism preserves the order valuation.
Restriction defines an action on $k[[T]]$ by $k$-algebra
automorphisms, and this action is faithful if and only if the
original action is faithful. If $G$ is a $p$-group, the restricted
action factors through $J_p$, and its Laurent extension recovers
the given action. If in addition $G$ is finite and discrete and
the action is faithful, this gives a closed topological embedding
$G\hookrightarrow J_p$.
\end{lemma}

\begin{proof}
The power-series ring inside $k((T))$ is precisely the set of
nonnegative-order elements. A valuation-preserving automorphism
and its inverse therefore preserve that ring. Their restrictions
are inverse algebra automorphisms, and restricting a composite
is the composite of the restrictions. Thus restriction gives a
group action, agreeing with the original action on every power
series.
Every Laurent series is a fraction of power series: multiplying
it by a sufficiently large nonnegative power of $T$ removes all
negative exponents. Consequently two field automorphisms which
agree on power series agree on every Laurent series, by applying
them to such a fraction. This proves preservation and reflection
of faithfulness, as well as recovery of the original action after
extending the restriction to the fraction field.
For a $p$-group the restricted automorphisms are normalized by
Lemma~\ref{lem:pgroup-power-series-normalization}, so their action
factors through $J_p$. That lemma also supplies the closed
embedding assertion for a faithful finite discrete action.
\end{proof}
